\section{Introduction}\label{sec:introduction}
Event-driven classifiers process not only where an event occurs but when it occurs. Retiming can change a prediction while preserving event identities, spatial coordinates, polarity and total count. Surrogate-gradient learning makes spiking dynamics trainable \cite{neftci2019}, while spike-retiming attacks make timing an explicit adversarial input dimension \cite{yu2026}. An evaluation must therefore specify both the admissible temporal changes and the computation available to find them.

We investigate end-to-end temporal vulnerability in an implemented hybrid classifier, emphasizing how computational matching changes attack comparisons and why results from separate temporal representations require separate interpretation. Our primary target, QSNN-v3, maps events through a classical Conv/LIF frontend, learned angles, an eight-qubit circuit and an affine probability readout. The central question is whether conclusions about its observed attack outcomes change with the definition of matched computation. Candidate count, forward-pass count and elapsed time measure different resources, particularly when only one attack uses gradients.

On the same 100 joint clean-correct validation examples, timestamp PGD causes three hybrid failures at each tested radius. TEMP-DRIFT-v2, a derivative-free margin search, causes 3--5 depending on radius and computational control. The largest numerical difference occurs under elapsed-time matching; no paired test establishes superiority. A separate four-model classical packet-grid study supplies context: similar clean accuracy need not imply similar attack outcomes. Differences in victims, populations, checkpoints and budget units prevent a causal representation ablation or a pooled robustness ranking across these studies.

The evaluation concerns the implemented composition. Quantum-classifier robustness can depend on encoding, model structure and adversarial access \cite{wendlinger2024comparative,gong2024randomized,dowling2026guarantees}; here the perturbation enters before the classical frontend. The following concrete outputs organize the paper:
\begin{enumerate}
\item \textbf{An explicit hybrid attack interface.} We specify temporal aggregation, Conv/LIF features and the sole quantum-probability readout, with preserved implementation evidence and a readout-specific margin observation. Missing artifacts limit full reproduction.
\item \textbf{Paired computational comparisons.} We report success counts, overlaps, exact paired tests and resource accounting for candidate, forward and elapsed-time controls on a frozen timestamp population.
\item \textbf{Classical interpretation controls.} We compare four architectures under a common packet-retiming recipe to delimit what clean accuracy and one optimizer establish.
\end{enumerate}
Source labels [R1]--[R10] map to preserved repository reports and code in the supplement. The supplement also retains development studies, full numerical tables and the conceptual literature comparison.
